\chapter{Experiments}\label{chap:experiments}

\section{Experiments summary}\label{sec:experiments-summary}

This chapter presents a structured evaluation of classical and neural relay strategies under progressively more challenging communication scenarios. 
To improve clarity and avoid repetition, all experimental configurations are first summarised in a single table.
Detailed numerical results, tables, and figures are provided exclusively in the appendices.
Each experiment subsection below therefore focuses solely on the key conclusions.

\begin{longtable}{p{0.6cm} p{2.2cm} p{2.1cm} p{3.6cm} p{3.9cm}}
\caption{Summary of experimental configurations (results reported in appendices)}
\label{tab:experiment-summary}\\
\hline
\textbf{Exp.} &
\textbf{Topology} &
\textbf{Modulation} &
\textbf{Channel / Equalization} &
\textbf{Focus} \\
\hline
\endfirsthead

\hline
\textbf{Exp.} &
\textbf{Topology} &
\textbf{Modulation} &
\textbf{Channel / Equalization} &
\textbf{Focus} \\
\hline
\endhead

E1 & SISO & BPSK &
AWGN, Rayleigh, Rician &
Channel model validation \\
\hline

E2 & SISO & BPSK &
AWGN, Rayleigh, Rician &
Baseline relay comparison \\
\hline

E3 & MIMO $2\times2$ & BPSK &
Rayleigh + ZF / MMSE / SIC &
Relay vs.\ MIMO equalization \\
\hline

E4 & SISO, MIMO $2\times2$ & BPSK &
AWGN, Rayleigh, Rician + EQs &
Parameter normalization \\
\hline

E5 & SISO & QPSK, 16-QAM &
AWGN, Rayleigh &
Higher-order modulation \\
\hline

E6 & SISO & 16-QAM &
AWGN &
CSI injection, input normalization \\
\hline

E7 & SISO & 16-QAM &
AWGN &
Joint 2D (16-class) classification \\
\hline

E8 & SISO & Learned ($M=16$) &
Rayleigh + ZF &
End-to-end autoencoder \\

\hline
\end{longtable}

\section{Channel Model Validation}
\label{sec:channel-model-validation}

The simulated BER curves closely match closed-form theoretical expressions for AWGN, Rayleigh, and Rician fading channels.
This confirms the correctness of the channel models and simulation framework.

\textbf{Relevant material:}
Validation tables and comparison figures are provided in
Appendix~\ref{chap:app-validation}, Figures~\ref{fig:app_awgn_validation}--
\ref{fig:app_rician_validation} and Tables~\ref{tab:app_channel_validation}.
\section{SISO BPSK Performance}
\label{sec:siso-bpsk}

At low SNR, learning-based relays outperform AF and, on selected channels, DF.
Under Rayleigh fading and universally at medium-to-high SNR, classical DF achieves equal or lower BER than all neural relays despite having no trainable parameters.

\textbf{Relevant material:}
BER tables and confidence-interval plots for AWGN, Rayleigh, and Rician channels
are provided in Appendix~\ref{chap:app-siso},
Figures~\ref{fig:app_siso_awgn}--\ref{fig:app_siso_rician}
and Tables~\ref{tab:app_siso_baseline}.
\section{MIMO $2\times2$ BPSK Performance}
\label{sec:mimo-bpsk}

The classical MIMO equalization hierarchy (ZF $<$ MMSE $<$ SIC) is preserved across all relay strategies.
While neural relays retain a low-SNR advantage under ZF and MMSE, DF combined with SIC yields the lowest BER overall.

\textbf{Relevant material:}
MIMO BER curves and tables for ZF, MMSE, and SIC equalization are reported in
Appendix~\ref{chap:app-mimo},
Figures~\ref{fig:app_mimo_zf}--\ref{fig:app_mimo_sic}
and Tables~\ref{tab:app_mimo_results}.
\section{Parameter Normalization and Complexity Trade-off}
\label{sec:parameter-complexity}

After equalizing parameter budgets, performance differences across architectures narrow substantially.
Very small MLPs remain competitive, while overly large models exhibit overfitting, indicating that parameter count is the dominant factor.

\textbf{Relevant material:}
Normalized-parameter BER tables and complexity trade-off figures are provided in
Appendix~\ref{chap:app-complexity},
Figures~\ref{fig:app_norm_all}--\ref{fig:app_complexity_tradeoff}
and Tables~\ref{tab:app_norm_results}.
\section{Higher-Order Modulation Scalability}
\label{sec:modulation-scalability}

QPSK mirrors BPSK performance due to I/Q independence.
For 16-QAM, bounded activations induce a structural BER floor, which is eliminated via constellation-aware activation scaling.

\textbf{Relevant material:}
All modulation-dependent BER curves and activation comparisons are provided in
Appendix~\ref{chap:app-modulation},
Figures~\ref{fig:app_qam16_awgn}--\ref{fig:app_activation_comparison}
and Tables~\ref{tab:app_modulation}.
\section{CSI Injection and Input Normalization}
\label{sec:csi-injection}

Explicit CSI injection improves performance for phase-only modulations but degrades performance for amplitude-based constellations.
Input normalization proves consistently beneficial.

\textbf{Relevant material:}
Complete CSI ablation results are reported in
Appendix~\ref{chap:app-csi},
Figures~\ref{fig:app_csi_ber} and Tables~\ref{tab:app_csi}.
\section{Joint 16-Class 2D Classification}
\label{sec:joint-classification}

Joint 2D classification eliminates the BER floor caused by per-axis I/Q splitting.
Several neural relays match classical DF at high SNR, confirming the limitation was architectural rather than fundamental.

\textbf{Relevant material:}
16-class vs.\ 4-class comparisons are provided in
Appendix~\ref{chap:app-16class},
Figures~\ref{fig:app_16class_ber}--\ref{fig:app_16class_heatmap}
and Tables~\ref{tab:app_16class}.
\section{End-to-End Joint Optimization}
\label{sec:end-to-end}

The end-to-end autoencoder underperforms both classical 16-QAM and modular DF relaying across all SNR points.
This confirms that unconstrained end-to-end learning is inferior to modular designs with embedded signal-processing structure.

\textbf{Relevant material:}
All end-to-end BER curves, learned constellations, and loss plots are provided in
Appendix~\ref{chap:app-e2e},
Figures~\ref{fig:app_e2e_ber}--\ref{fig:app_e2e_constellation}.
